\section{量产视角：技术、质量、成本和责任要一起算}

前面讲了技术路线和组织简化，本章补上量产视角。刘先明多次强调，他更喜欢从真实问题里找下一阶段要做的事情。自动驾驶负责人不只是研究负责人，还要面对产品、业务、质量、硬件、成本和用户安全。这个视角解释了为什么车企不能简单照搬研究团队的节奏。

\subsection{量产不是研究的尾声，而是问题来源}

研究团队常常先定义任务，再找数据和模型；量产团队面对的是真实用户每天遇到的问题。车在路上出现的每个接管、每个不舒适动作、每个误判、每个用户投诉，都会变成下一轮模型和工程要解决的问题。量产不是研究结束后的部署环节，而是持续产生研究问题的入口。

\begin{knowledgebox}{量产问题如何反哺模型}
\begin{tabularx}{\textwidth}{p{0.22\textwidth}p{0.34\textwidth}X}
\toprule
量产信号 & 进入模型前要做什么 & 反哺价值 \\
\midrule
接管 & 定位接管原因、切出片段、标注上下文 & 找到模型盲区和规则上限。 \\
用户不适 & 关联驾驶动作、速度、车距和环境 & 优化舒适性和风格。 \\
事故/风险 & 做责任归因、复盘感知/预测/规划链路 & 形成安全样本和评测集。 \\
硬件约束 & 记录算力、传感器、延迟和成本限制 & 指导蒸馏、量化和剪枝。 \\
\bottomrule
\end{tabularx}
\end{knowledgebox}

\subsection{为什么“简单即美”不是粗暴删除}

简单不是把所有保护都删掉，而是把重复、过时、阻碍 scaling 的结构删掉。车企仍然必须保留安全冗余、质量验证和上线流程。刘先明强调的简单，是让研发工序更短、目标更少、问题暴露更快，而不是牺牲安全边界。

\begin{warningbox}{简化不等于冒进}
如果为了追求端到端而忽略安全、测试和量产质量，简化会变成冒进。真正的简化，是减少无效复杂度，同时保留必要安全约束。
\end{warningbox}

\subsection{本章小结}

Physical AI 的量产路线要求技术和责任一起前进。模型要更大、数据要更多，但上线必须经受安全、质量、成本和用户体验检验。

